\section*{CHƯƠNG 3. PHƯƠNG PHÁP ĐỀ XUẤT}
\addcontentsline{toc}{section}{\numberline{}CHƯƠNG 3. PHƯƠNG PHÁP ĐỀ XUẤT}
\setcounter{section}{3}
\setcounter{subsection}{0}
\setcounter{figure}{0}
\setcounter{table}{0}

\subsection*{3.1. Quy trình tổng quát của phương pháp nghiên cứu đề xuất}
\addcontentsline{toc}{subsection}{\numberline{}3.1. Quy trình tổng quát của phương pháp nghiên cứu đề xuất}

Để giải quyết bài toán nhận dạng danh tính khách hàng bị biến đổi diện mạo do phẫu thuật thẩm mỹ (PTTM) hoặc lão hóa, nghiên cứu này đề xuất phương pháp trích xuất đặc trưng phân vùng thích ứng, kết hợp ánh xạ không gian metric qua lớp chiếu tuyến tính được huấn luyện bằng hàm mất mát Triplet Loss. 

Quy trình nghiên cứu tổng quát được chia làm hai giai đoạn chính:
\begin{enumerate}
    \item \textbf{Giai đoạn huấn luyện offline (Offline Training):} 
    Sử dụng tập dữ liệu ảnh khuôn mặt PTTM để thực hiện phân tách vùng $\rightarrow$ Trích xuất đặc trưng thô từ mô hình nền tảng Facenet512 $\rightarrow$ Huấn luyện lớp chiếu mạng nơ-ron bằng Triplet Loss để tối ưu hóa ma trận trọng số $\bm{W}$ và bias $\bm{b}$ cho từng vùng ảnh độc lập.
    \item \textbf{Giai đoạn đối sánh online (Online Matching):} 
    Ảnh đầu vào từ thiết bị được phân vùng và chiếu qua ma trận trọng số đã học để thu được các vector đặc trưng thích ứng $\rightarrow$ Thực hiện thuật toán đối sánh dung hợp trọng số sinh học để đưa ra quyết định nhận dạng danh tính cuối cùng.
\end{enumerate}

\newpage
\subsection*{3.2. Phương pháp tiền xử lý và phân tách vùng ảnh khuôn mặt}
\addcontentsline{toc}{subsection}{\numberline{}3.2. Phương pháp tiền xử lý và phân tách vùng ảnh khuôn mặt}

Các mô hình nhận dạng khuôn mặt sâu quy mô lớn hiện nay như DeepFace \cite{deepface_framework} và ArcFace \cite{arcface} đều tiếp cận theo phương pháp học đặc trưng toàn cục (holistic feature learning). Các mô hình này tiếp nhận toàn bộ hình ảnh khuôn mặt đã căn chỉnh làm đầu vào và ánh xạ chúng thành một vector đặc trưng duy nhất đại diện cho tổng thể khuôn mặt. Tuy nhiên, các nghiên cứu thực nghiệm của Singh và cộng sự \cite{singh2010database} cũng như Rathgeb và cộng sự \cite{rathgeb2020obstacle} đã chỉ ra rằng phương pháp toàn cục này rất nhạy cảm với các biến đổi diện mạo cục bộ: chỉ một thay đổi nhỏ do phẫu thuật thẩm mỹ ở cằm hoặc mũi sẽ làm méo mó các giá trị kích hoạt lan truyền trong mạng nơ-ron, dẫn đến sai lệch lớn trên toàn bộ vector đặc trưng đầu ra. 

Để giải quyết điểm yếu này, một số nghiên cứu tiền đề đã chuyển hướng sang trích xuất đặc trưng cục bộ (local feature learning), tiêu biểu là phương pháp phân vùng LBP kết hợp của Bhatt và cộng sự \cite{iiitd_plastic_surgery}. Kế thừa tư duy phân vùng đó, phương pháp đề xuất của nghiên cứu này thực hiện phân tách khuôn mặt thành ba phân vùng sinh học độc lập theo chiều dọc để cô lập các biến dạng cơ học:
\begin{itemize}
    \item \textbf{Vùng Thượng (Upper Region):} Bao gồm trán, lông mày và hốc mắt. Đây là vùng có cấu trúc xương cố định nhất, ít bị can thiệp bởi các ca PTTM phổ biến.
    \item \textbf{Vùng Trung (Mid Region):} Bao gồm sống mũi, hai bên má và tai. Vùng này chịu ảnh hưởng trực tiếp từ phẫu thuật nâng mũi (Rhinoplasty) hoặc tạo hình gò má.
    \item \textbf{Vùng Hạ (Lower Region):} Bao gồm nhân trung, miệng, môi, hàm và cằm. Vùng này chịu ảnh hưởng mạnh nhất từ các phẫu thuật gọt hàm, độn cằm, cắt môi.
\end{itemize}

Để thực hiện chia cắt không để lại các vết cắt sắc cạnh gây nhiễu cho mạng tích chập, nghiên cứu áp dụng kỹ thuật \textbf{Mặt nạ làm mờ biên (Blur-Masking)}:
\begin{enumerate}
    \item Xác định các tọa độ mốc giới hạn phân chia dựa trên tỷ lệ chiều cao khuôn mặt sau khi căn chỉnh chính diện.
    \item Tạo các mặt nạ nhị phân tương ứng với từng vùng Thượng, Trung, Hạ.
    \item Áp dụng bộ lọc Gaussian Blur lên biên của mặt nạ để làm mịn vùng chuyển tiếp giữa các lát cắt, đảm bảo thông tin biên không bị mất đột ngột khi đưa vào mô hình học sâu trích xuất đặc trưng.
\end{enumerate}

\subsection*{3.3. Thiết kế mạng ánh xạ đặc trưng thích ứng (Feature Projector)}
\addcontentsline{toc}{subsection}{\numberline{}3.3. Thiết kế mạng ánh xạ đặc trưng thích ứng (Feature Projector)}

Lớp chiếu đặc trưng đề xuất (Projector) được thiết kế như một lớp liên kết đầy đủ (Fully Connected Layer) đơn giản đặt ngay sau lớp đầu ra của mô hình nền tảng Facenet512.

Với mỗi vector đặc trưng thô $\bm{x} \in \mathbb{R}^{512}$ trích xuất từ một vùng ảnh bởi Facenet512, lớp chiếu thực hiện phép biến đổi tuyến tính để ánh xạ sang vector đặc trưng mới $\bm{y} \in \mathbb{R}^{512}$ theo công thức:
$$\bm{z} = \bm{x} \bm{W} + \bm{b}$$
$$\bm{y} = \frac{\bm{z}}{\|\bm{z}\|_2}$$
Trong đó:
\begin{itemize}
    \item $\bm{W} \in \mathbb{R}^{512 \times 512}$ là ma trận trọng số cần tối ưu hóa.
    \item $\bm{b} \in \mathbb{R}^{512}$ là vector bias độ lệch.
    \item Phép chia cho chuẩn bậc hai $\|\bm{z}\|_2$ (L2 Normalization) để đưa vector đầu ra về mặt cầu đơn vị, phục vụ tính toán khoảng cách Cosine.
\end{itemize}

\subsection*{3.4. Quy trình huấn luyện tối ưu hóa ma trận chiếu bằng Triplet Loss}
\addcontentsline{toc}{subsection}{\numberline{}3.4. Quy trình huấn luyện tối ưu hóa ma trận chiếu bằng Triplet Loss}

Ma trận chiếu $\bm{W}$ và bias $\bm{b}$ của lớp chiếu tuyến tính được huấn luyện tối ưu hóa thông qua cơ chế học khoảng cách đặc trưng (Metric Learning). Trong đó, **hàm mất mát Triplet Loss** là một thuật toán tối ưu hóa khoảng cách đặc trưng kinh điển được tác giả tìm hiểu, kế thừa từ nghiên cứu gốc của Schroff và các cộng sự (mô hình FaceNet) và chủ động lập trình cài đặt để huấn luyện lớp chiếu thích ứng tự thiết kế cho từng vùng khuôn mặt.

\subsubsection*{3.4.1. Thiết lập bộ ba huấn luyện (Triplet Selection)}
Với mỗi bước huấn luyện, hệ thống xây dựng các bộ ba mẫu (triplets) $(A, P, N)$ từ tập dữ liệu phẫu thuật thẩm mỹ:
\begin{itemize}
    \item \textbf{Anchor ($A$):} Ảnh khuôn mặt gốc trước PTTM của chủ thể $i$.
    \item \textbf{Positive ($P$):} Ảnh khuôn mặt sau PTTM của cùng chủ thể $i$.
    \item \textbf{Negative ($N$):} Ảnh khuôn mặt của chủ thể khác $j$ ($j \neq i$).
\end{itemize}

\subsubsection*{3.4.2. Hàm mục tiêu tối ưu hóa}
Hàm mất mát Triplet Loss được định nghĩa để cực tiểu hóa khoảng cách giữa hai ảnh của cùng một người đồng thời cực đại hóa khoảng cách đối với người khác:
$$\mathcal{L}(\bm{W}, \bm{b}) = \sum_{i=1}^{M} \max \left( \|f(A_i) - f(P_i)\|_2^2 - \|f(A_i) - f(N_i)\|_2^2 + \alpha, 0 \right)$$
Trong đó $f(x)$ là kết quả vector đầu ra của lớp chiếu đã chuẩn hóa L2 của ảnh $x$, và $\alpha = 0.2$ là biên an toàn (margin). Lan truyền ngược (Backpropagation) được áp dụng để cập nhật ma trận $\bm{W}$ và bias $\bm{b}$ sao cho hàm mất mát $\mathcal{L}$ tiệm cận về 0.

\subsection*{3.5. Thuật toán đối sánh dung hợp trọng số sinh học đề xuất (Weighted Bio-fusion)}
\addcontentsline{toc}{subsection}{\numberline{}3.5. Thuật toán đối sánh dung hợp trọng số sinh học đề xuất (Weighted Bio-fusion)}

Quy trình đối sánh dung hợp trọng số sinh học là thuật toán do tác giả tự đề xuất thiết kế và lập trình cài đặt dựa trên tổ hợp tuyến tính các khoảng cách Cosine cục bộ. Thuật toán này tích hợp bộ trọng số sinh học tối ưu đã được chứng minh qua thực nghiệm tìm kiếm lưới ở Chương 2.

Trong pha đối sánh thực tế, thay vì so khớp một khoảng cách duy nhất, hệ thống tính toán khoảng cách Cosine giữa các phân vùng ảnh tương ứng và dung hợp chúng theo công thức trọng số sinh học:
$$D_{\text{fusion}} = w_{\text{upper}} \cdot D_{\text{Cosine}}(A_{\text{upper}}, B_{\text{upper}}) + w_{\text{mid}} \cdot D_{\text{Cosine}}(A_{\text{mid}}, B_{\text{mid}}) + w_{\text{lower}} \cdot D_{\text{Cosine}}(A_{\text{lower}}, B_{\text{lower}})$$
Trong đó, các hệ số trọng số được gán dựa trên mức độ ổn định cấu trúc sinh học trước PTTM:
$$w_{\text{upper}} = 0.5; \quad w_{\text{mid}} = 0.3; \quad w_{\text{lower}} = 0.2$$
Nếu khoảng cách dung hợp $D_{\text{fusion}} < \theta$ (với ngưỡng đối sánh $\theta = 0.30$), hệ thống đưa ra kết luận hai khuôn mặt thuộc về cùng một cá nhân.

\subsection*{3.6. Mô hình hệ thống thử nghiệm phục vụ thực nghiệm (Demo System)}
\addcontentsline{toc}{subsection}{\numberline{}3.6. Mô hình hệ thống thử nghiệm phục vụ thực nghiệm (Demo System)}

Để thực nghiệm và chứng minh tính ứng dụng của phương pháp nghiên cứu đề xuất, một mô hình ứng dụng thử nghiệm (Demo) được xây dựng dưới cấu trúc kiến trúc 3 lớp:
\begin{itemize}
    \item \textbf{Giao diện Demo (Next.js):} Mô phỏng các chức năng đăng ký chân dung thành viên và quét camera nhận diện.
    \item \textbf{Dịch vụ xử lý (FastAPI):} Đóng gói mô hình đề xuất và thực hiện tính toán đặc trưng.
    \item \textbf{Hệ lưu trữ đặc trưng (PostgreSQL + pgvector):} Lưu trữ thông tin định danh mẫu và hỗ trợ tính toán khoảng cách đối sánh dung hợp trực tiếp ở mức CSDL bằng ngôn ngữ truy vấn SQL.
\end{itemize}
